\subsection{局部紧群的酉对偶}

设 G 是局部紧拓扑群。给 G 赋予左哈尔测度 \(\mu\)。对 \(p \in [1, +\infty]\)，以 \(\mathcal{L}^{p}(\mathrm{G})\)（分别以 \(\mathrm{L}^{p}(\mathrm{G})\)）表示空间 \(\mathcal{L}_{\mathbf{C}}^{p}(\mathrm{G}, \mu)\)（分别表示空间 \(\mathrm{L}_{\mathbf{C}}^{p}(\mathrm{G}, \mu)\)）。将空间 \(\mathcal{C}_{b}(\mathrm{G})\) 与其在 \(\mathrm{L}^{\infty}(\mathrm{G})\) 中的像等同起来。

以 \(\Delta\) 表示 G 的模函数。回忆 \(L^1(G)\) 是对合巴拿赫代数，其对合由映射 \(f \mapsto f^*\) 通过取商诱导，其中 \(f^*(y) = \Delta^{-1}(y)\overline{f(y^{-1})}\) 对所有 \(f \in \mathcal{L}^1(G)\) 和 \(y \in G\) 成立（参见 I, p. 99, 例 4）。根据 INT, VIII, p. 172, §4, no. 7, 命题 20，巴拿赫代数 \(L^1(G)\) 具有逼近单位元。

设 \(\varrho\) 是 G 在复希尔伯特空间 E 上的酉表示。映射 \(f\mapsto\varrho(f)\) 是从 L \(^{1}\) (G) 到 \(\mathcal{L}\) (E) 的连续对合代数态射（V, p. 401, 引理 1）；它是 L \(^{1}\) (G) 在 E 上的非退化表示（INT, VIII, p. 139, § 2, no. 7, 命题 10, (i)）。

命题 12. --- 设 E 是复希尔伯特空间，\(\widetilde{\pi}\) 是从 L \(^{1}\) (G) 到 \(\mathcal{L}\) (E) 的对合代数态射。若表示 \(\widetilde{\pi}\) 非退化，则存在 G 在 E 上的唯一酉表示 \(\pi\)，使得 \(\widetilde{\pi}(f)=\pi(f)\) 对所有 \(f\in\mathrm{L}^{1}(\mathrm{G})\) 成立。

唯一性由 INT, VIII, p. 139, § 2, n° 7, 引理 4 的推论 3 得出。

证明 \(\pi\) 的存在性。由 I, p. 104, 命题 2，有 \(\| \widetilde{\pi}\| \leqslant 1\)。令 F 为 E 中由形如 \(\widetilde{\pi}(f)x\) 的向量生成的子空间，其中 \(f\) 属于 \(\mathrm{L}^1 (\mathrm{G})\)，\(x\) 属于 E；由于表示 \(\widetilde{\pi}\) 非退化，F 在 E 中稠密。

对 e 的每个紧邻域 V，令 \(\varphi_{V}\) 为支集包含于 V 且满足 \(\int\varphi_{V}\mu=1\) 的连续正函数。令 B 为 e 的紧邻域滤子的一个基。

设 \(g\) 属于 \(G\)，设 \(f\) 属于 \(L^1(G)\)。有 \((\varphi_V * \varepsilon_g) * f \to \varepsilon_g * f\) 在 \(L^1(G)\) 中沿 \(\mathfrak{B}\) 的截面滤子收敛（INT, VIII, p. 172, § 4, no. 7, 命题 20），所以 \(\widetilde{\pi}(\varphi_V * \varepsilon_g)\widetilde{\pi}(f)\) 收敛到 \(\widetilde{\pi}(\varepsilon_g * f)\)（在 \(\mathcal{L}(E)\) 中）。这意味着 \(\widetilde{\pi}(\varphi_V * \varepsilon_g)\) 沿 \(\mathfrak{B}\) 的截面滤子逐点收敛于映射 \(\pi(g)\)，从 \(F\) 到 \(F\)；由于 \(\| \widetilde{\pi}(\varphi_V * \varepsilon_g) \| \leqslant 1\) 对每个 \(V\)（它是 \(e\) 的紧邻域）成立，该映射由 EVT, III, p. 26, 推论 3 可知是连续的。因此存在唯一的 \(E\) 的自同态，它通过对子空间的传递诱导出 \(\pi(g)\)。仍以 \(\pi(g)\) 表示该自同态；有 \(\| \pi(g) \| \leqslant 1\)。

设 \(f \in \mathrm{L}^1(\mathrm{G})\)。由于 \(\widetilde{\pi}(\varphi_{\mathrm{V}} * \varepsilon_g) \widetilde{\pi}(f)\) 收敛到 \(\widetilde{\pi}(\varepsilon_g * f)\)，有 \(\pi(g) \widetilde{\pi}(f) = \widetilde{\pi}(\varepsilon_g * f)\)。

当 \(g = e\) 时，该关系表明 \(\pi(e)\) 在 \(F\) 上是恒等映射，因而 \(\pi(e) = 1_{E}\)。设 \(g_1\) 和 \(g_2\) 属于 \(G\)。有

\[
\pi (g _ {1}) \pi (g _ {2}) \widetilde {\pi} (f) = \pi (g _ {1}) \widetilde {\pi} (\varepsilon_ {g _ {2}} * f) = \widetilde {\pi} (\varepsilon_ {g _ {1} g _ {2}} * f) = \pi (g _ {1} g _ {2}) \widetilde {\pi} (f),
\]

所以 \(\pi(g_1)\pi(g_2) = \pi(g_1g_2)\) 在 F 上成立，因而在 E 上也成立。这证明了映射 \(g \mapsto \pi(g)\) 是 G 在 E 上的表示。由于 \(\|\pi(g)\| \leqslant 1\) 且 \(\|\pi(g)^{-1}\| = \|\pi(g^{-1})\| \leqslant 1\)，E 的自同态 \(\pi(g)\) 是等距的，因而是酉的（EVT, V, p. 40, 命题 3）。

设 \(f \in \mathrm{L}^{1}(\mathrm{G})\) 且 \(x \in E\)。映射 \(g \mapsto \varepsilon_{g} * f\) 从 G 到 \(\mathrm{L}^{1}(\mathrm{G})\) 连续（INT, VIII, p. 136, § 2, no. 5 及 p. 144, § 3, no. 2, 公式 (5)），且 \(\widetilde{\pi}\) 连续，因此映射 \(g \mapsto \widetilde{\pi}(\varepsilon_{g} * f)x = \pi(g)\widetilde{\pi}(f)x\) 在 \(g = e\) 处连续。由于 F 在 E 中稠密，表示 \(\pi\) 是酉表示（V, p. 380, 引理 4）。

设 \(f_1\) 和 \(f_2\) 属于 \(L^1(G)\)。根据 INT, VIII, p. 127, §1, no. 5, 命题 7，有关系

\[
f _ {1} * f _ {2} = \int_ {\mathrm{G}} f _ {1} (g) (\varepsilon_ {g} * f _ {2}) d \mu (g)
\]

在 L \(^{1}\) (G) 中，从而

\[
\begin{array}{l} \widetilde {\pi} (f _ {1}) \widetilde {\pi} (f _ {2}) = \widetilde {\pi} (f _ {1} * f _ {2}) = \int_ {\mathrm{G}} f _ {1} (g) \widetilde {\pi} (\varepsilon_ {g} * f _ {2}) d \mu (g) \\ \qquad = \int_ {\mathrm{G}} f _ {1} (g) \pi (g) \widetilde {\pi} (f _ {2}) d \mu (g) = \Big (\int_ {\mathrm{G}} f _ {1} (g) \pi (g) d \mu (g) \Big) \widetilde {\pi} (f _ {2}), \end{array}
\]

利用 INT, VI, p. 9, § 1, no. 1, 命题 1，可得

\[
\widetilde {\pi} (f) = \int_ {\mathrm{G}} f (g) \pi (g) d \mu (g) = \pi (f)
\]

对所有 \(f \in \mathrm{L}^{1}(\mathrm{G})\) 都成立。命题得证。

命题 13. --- 设 \(\varphi\in\mathrm{L}^{\infty}(\mathrm{G})\)。以 \(\lambda_{\varphi}\) 表示由 \(f\mapsto\langle f,\varphi\rangle\) 定义在 \(\mathrm{L}^{1}(\mathrm{G})\) 上的连续线性泛函。则 \(\varphi\) 是 G 上某个正定函数的等价类，当且仅当 \(\lambda_{\varphi}\) 是 \(\mathrm{L}^{1}(\mathrm{G})\) 上的连续正线性泛函。

设 \(\lambda_{\varphi}\) 是 \(L^{1}(G)\) 上的连续正线性泛函。令 \((E, x, \widetilde{\pi})\) 为 \(\lambda\) 的循环希尔伯特实现（V, p. 438, 命题 5）。

令 \(\pi\) 为 G 在 E 上的酉表示，使得 \(\pi(f) = \widetilde{\pi}(f)\) 对所有 \(f \in \mathrm{L}^1(\mathrm{G})\) 成立（命题 12）。于是 \(\lambda_{\varphi}(f) = \langle x | \pi(f)x \rangle\) 对所有 \(f \in \mathrm{L}^1(\mathrm{G})\) 成立。

由于

\[
\pi (f) = \int_ {\mathrm{G}} f (g) \pi (g) d \mu (g)
\]

对所有 \(f\in \mathrm{L}^1 (\mathrm{G})\) 都成立，从而

\[
\int_ {\mathrm{G}} f (g) \varphi (g) d \mu (g) = \lambda_ {\varphi} (f) = \langle x | \pi (f) x \rangle = \int_ {\mathrm{G}} f (g) \langle x | \pi (g) x \rangle d \mu (g)
\]

对所有 \(f \in \mathrm{L}^{1}(\mathrm{G})\) 都成立。因此，\(\varphi\) 是 \(\mathrm{L}^{\infty}(\mathrm{G})\) 中由 G 上的函数 \(g \mapsto \langle x | \pi(g)x \rangle\) 所定义的等价类，而该函数是 G 上的正定函数（V, p. 443, 命题 9）。

反之，设 \(\varphi\in\mathrm{Pos}(\mathrm{G})\)。令 \(f\in\mathcal{K}(\mathrm{G})\)。于是有

\[
\begin{array}{c} \lambda_ {\varphi} (f ^ {*} * f) = \int_ {G} \varphi (y) \Big (\int_ {G} \Delta (z) ^ {- 1} \overline {{f (z ^ {- 1})}} f (z ^ {- 1} y) d \mu (z) \Big) d \mu (y) \\ = \int_ {G} \varphi (y) \Big (\int_ {G} \overline {{f (z)}} f (z y) d \mu (z) \Big) d \mu (y) \end{array}
\]

G × G 上由 \((z, y) \mapsto \varphi(y)\overline{f(z)}f(zy)\) 定义的连续函数是有界的，并且

\[
\begin{array}{c} \int_ {\mathrm{G}} ^ {*} | \varphi (y) | \Big (\int_ {\mathrm{G}} ^ {*} | \overline {{f (z)}} f (z y) | d \mu (z) \Big) d \mu (y) \\ \leqslant \varphi (e) \int_ {\mathrm{G}} ^ {*} \int_ {\mathrm{G}} ^ {*} | f (z) | | f (z y) | d \mu (z) d \mu (y) \\ = \varphi (e) \Big (\int_ {\mathrm{G}} ^ {*} | f (z) | d \mu (z) \Big) ^ {2} \end{array}
\]

由 INT, V, p. 94, § 8, no. 3, 命题 8，利用勒贝格–富比尼定理（INT, V, p. 96, § 8, no. 4, 定理 1），可推出

\[
\begin{array}{c} \lambda_ {\varphi} (f ^ {*} * f) = \int_ {\mathrm{G}} \overline {{f (z)}} \Bigl (\int_ {\mathrm{G}} \varphi (y) f (z y) d \mu (y) \Bigr) d \mu (z) \\ = \int_ {\mathrm{G} \times \mathrm{G}} \overline {{f (z)}} f (y) \varphi (z ^ {- 1} y) d (\mu \otimes \mu) (z, y). \end{array}
\]

由 V, p. 432, 定理 1, (i)，可知 \(\lambda_{\varphi}(f^{*} * f) \geqslant 0\)，这是将其应用于 \(\mu\) 在 \(f\) 的支集上诱导的测度的结果（INT, IV, p. 186, § 5, n° 7, 定义 4）。由连续性可推出，\(\lambda_{\varphi}(f^{*} * f) \geqslant 0\) 对所有 \(f \in \mathrm{L}^1(\mathrm{G})\) 成立。

空间 \(\mathrm{L}^{\infty}(\mathrm{G})\) 和 \(\mathcal{C}_{b}(\mathrm{G})\) 赋予其巴拿赫空间拓扑，并以 \(f \mapsto \|f\|_{\infty}\) 表示范数。这里将 \(\mathrm{L}^{\infty}(\mathrm{G})\)、\(\mathcal{C}_{b}(\mathrm{G})\) 或 \(\operatorname{Pos}(\mathrm{G})\) 上由弱拓扑 \(\sigma(\mathrm{L}^{\infty}(\mathrm{G}), \mathrm{L}^{1}(\mathrm{G}))\) 诱导的拓扑称为弱拓扑。

由于 \(L^{\infty}(G)\) 与 \(L^{1}(G)\) 的对偶空间相等同（INT, V, p. 61, §5, no. 8, 定理 4），\(L^{\infty}(G)\) 的每个闭球对于弱拓扑都是紧的（EVT, III, p. 17, 推论 3）。

推论。--- 在赋予弱拓扑的 \(\mathrm{L}^{\infty}(\mathrm{G})\) 中，集合 \(\operatorname{Pos}(\mathrm{G})\) 是闭的，集合 \(\operatorname{Pos}_{\leqslant 1}(\mathrm{G})\) 是紧的。

第一个断言由命题得出，第二个断言则由 EVT, III, 同上得出。

一般而言，\(\mathrm{Pos}_{1}(\mathrm{G})\) 对弱拓扑并不紧，群 R 的例子说明了这一点（V, p. 497, 习题 10）。

命题 14. --- \(\mathrm{Pos}_{\leqslant 1}(\mathrm{G})\) 的极点集合等于 \(\mathrm{Pos}_1(\mathrm{G})\) 的极点集合与零函数的并。此外，\(\mathrm{Pos}_1(\mathrm{G})\) 的极点集合的闭凸包包含 \(\mathrm{Pos}_1(\mathrm{G})\)。

由命题 13 的推论，集合 \(\mathrm{Pos}_{\leqslant 1}(\mathrm{G})\) 是尖凸锥 \(\mathrm{Pos}(\mathrm{G})\) 的一个截面，该截面由规范 \(\varphi \mapsto \varphi(e)\) 定义（EVT, II, p. 61, 定义 3 及命题 4）。因此，它的极点是零函数以及元素 \(\varphi\)，它属于 \(\mathrm{Pos}_1(\mathrm{G})\) 且属于 \(\mathrm{Pos}(\mathrm{G})\) 的极端生成元（EVT, II, p. 62, 推论 1）。这些正是 \(\mathrm{Pos}_1(\mathrm{G})\) 的极点（EVT, II, p. 61, 命题 3）。第一个断言由此得出。

证明第二个断言。设 \(\varphi\in\mathrm{Pos}_{1}(\mathrm{G})\)。由于集合 \(\mathrm{Pos}_{\leqslant 1}(\mathrm{G})\) 对弱拓扑是紧的（命题 13 的推论），存在滤子 \(\mathfrak{F}\)，在 \(\mathrm{Pos}_{\leqslant 1}(\mathrm{G})\) 的极点集合的凸包上收敛于 \(\varphi\)（EVT, II, p. 59, 定理 1）。由于巴拿赫空间 \(\mathrm{L}^{\infty}(\mathrm{G})\) 的闭球对于弱拓扑是闭的，且 \(\|\varphi\|_{\infty}=1\)，有 \(\lim_{\psi,\mathfrak{F}}\|\psi\|_{\infty}=1\)（事实上，否则将存在实数 \(c<1\) 以及滤子 \(\mathfrak{G}\)，它在半径为 \(c\) 的 \(\mathrm{L}^{\infty}(\mathrm{G})\) 闭球上，且比 \(\mathfrak{F}\) 更细，这将意味着 \(\varphi\) 属于该闭球）。

根据对 \(\mathrm{Pos}_{\leqslant1}(\mathrm{G})\) 极点的描述，极点的凸包中的每个元 \(\psi\) 都是 G 上具有如下形式的 \(\mathrm{Pos}_{\leqslant1}(\mathrm{G})\) 上的正定函数

\[
\psi = \sum_ {i \in \mathrm{I}} t _ {i} \psi_ {i},
\]

其中 I 是有限集，\(\psi_{i}\) 是 \(\mathrm{Pos}_{1}(\mathrm{G})\) 的极点，对所有 \(i \in I\)，且 \(t_{i} \in [0,1]\)。有 \(\sum_{i} t_{i} = \psi(e) = \|\psi\|_{\infty} \leqslant 1\)（V, p. 446, 引理 4）。若 \(\psi \neq 0\)，函数

\[
\frac {\psi}{\| \psi \| _ {\infty}} = \sum_ {i \in \mathrm{I}} \frac {t _ {i}}{\psi (e)} \psi_ {i}
\]

因此属于 \(\mathrm{Pos}_{1}(\mathrm{G})\) 极点的凸包。由于 \(\lim_{\psi,\mathfrak{F}}\|\psi\|_{\infty}^{-1}\psi=\varphi\)，可知 \(\varphi\) 属于 \(\mathrm{Pos}_{1}(\mathrm{G})\) 极点的闭凸包。断言得证。

